\documentclass[conference]{IEEEtran}

\usepackage{cite}
\usepackage{amsmath,amssymb,amsfonts}
\usepackage{graphicx}
\usepackage{textcomp}
\usepackage{xcolor}
\usepackage{url}

\begin{document}

\title{Ultra-Nano-vLLM: A Reproducible Performance-Research Workflow for Lightweight LLM Inference}

\author{\IEEEauthorblockN{CHEN EN LO}
\IEEEauthorblockA{\textit{Computer Science Department} \\
\textit{National Yang Ming Chiao Tung University}\\
asciibase64.cs14@nycu.edu.tw}}

\maketitle

\begin{abstract}
Optimizing a large language model inference engine requires measurements that
separate a code change from workload variation, persistent accelerator state,
and reporting differences. Ultra-Nano-vLLM is a performance-research project
built around nano-vLLM. It provides a scalar benchmark lifecycle, validated
YAML experiment grids, fresh-process isolation, typed JSON reports, comparison
plots, optional profiling annotations, and file-selected implementations for
five replaceable engine roles. These facilities connect benchmarking,
controlled experimentation, profiling, focused optimization, and regression
comparison in one workflow while keeping the benchmark and experiment
entrypoints independent of concrete engine implementations. This paper
describes the workflow, its implementation boundaries, and the evaluation
protocol supplied by the repository.
\end{abstract}

\section{Introduction}

Large language model (LLM) inference systems coordinate model execution,
request scheduling, and key--value (KV) cache management under changing input
and output lengths. vLLM addresses KV-cache memory management through
PagedAttention and a block-based serving design~\cite{kwon2023vllm}.
nano-vLLM provides a lightweight implementation of this class of inference
engine~\cite{nanovllm}. Its compact implementation makes individual engine
components accessible, but an optimization study also needs a controlled way
to define workloads, preserve comparison conditions, collect measurements, and
separate successive runs.

Performance results are difficult to interpret when a benchmark changes more
than one relevant condition or when runtime state survives across
configurations. Reproducible inference benchmarking therefore requires an
explicit workload, consistent metrics, and comparable execution rules; the
broader need for comparability across inference systems also motivates
MLPerf Inference~\cite{reddi2019mlperf}. Within a CUDA-based LLM engine,
additional state can arise from model initialization, memory allocation,
compilation, CUDA Graph capture, and distributed process groups. A research
workflow must keep these effects outside the comparison or recreate them for
each scalar configuration.

Ultra-Nano-vLLM turns nano-vLLM into a controlled environment for performance
research. Its goal is not to extend the inference engine for its own sake, but
to support a repeatable loop: establish a baseline, isolate a factor through a
controlled experiment, profile the confirmed workload, apply a focused change,
and compare the resulting reports for improvement and regression. For example,
the repository can cross a set of named workloads with baseline and alternative
scheduler implementations while holding the remaining component selectors
fixed.

Three engineering requirements shape the workflow. First, every comparison
must resolve to explicit scalar configurations and reject incomplete or
ambiguous grids before GPU execution. Second, sequential configurations must
not share CUDA Graph pools, cached allocations, or communication state. Third,
the experiment entrypoint must be able to replace one scheduler, block manager,
attention implementation, sampler, or KV-cache storage kernel without learning
the semantics of that implementation.

The system addresses these requirements with separate benchmark, experiment,
reporting, and component-selection boundaries. A scalar benchmark owns the
measurement lifecycle. The experiment layer normalizes YAML input, validates
the requested grid, and runs each resolved configuration in a fresh spawned
process. Typed results return to the parent process for reporting and plotting.
An engine-component selector maps trusted local file stems to runtime-checked
factory contracts, allowing the same experiment machinery to transport
alternative implementations without branching on their roles.

This paper makes three contributions. First, it describes a reproducible
benchmark-to-optimization workflow for a lightweight LLM inference engine
(Sections~II and III). Second, it presents a validated experiment architecture
that combines scalar configuration, process isolation, typed reporting, and
dimension-specific comparison plots (Section~III). Third, it defines a
pluggable boundary for controlled component and kernel replacement together
with an evaluation protocol for scheduler comparisons (Sections~III and IV).

\section{Background}

\subsection{Inference Engine Foundation}

Ultra-Nano-vLLM is built on nano-vLLM, which describes itself as a lightweight
vLLM implementation and provides an offline generation interface
~\cite{nanovllm}. The upstream engine supplies the inference implementation;
Ultra-Nano-vLLM adds benchmark configuration, controlled experiment execution,
report persistence, plotting, profiling hooks, and pluggable implementation
selection. The resulting separation keeps performance-research concerns out of
the benchmarked component implementations.

The underlying memory organization uses fixed-size KV-cache blocks. In the
broader vLLM design, PagedAttention stores logical KV blocks in non-contiguous
physical memory and allocates additional blocks as a sequence grows
~\cite{kwon2023vllm}. Ultra-Nano-vLLM exposes peak used blocks and peak block
utilization as measurements so that latency and throughput changes can be
examined alongside cache pressure.

\subsection{Prefill and Decode}

Autoregressive LLM inference contains a prefill phase that processes prompt
tokens and a decode phase that repeatedly generates subsequent tokens.
Sarathi studies the different execution characteristics of these phases and
uses chunked prefills with decode-maximal batching~\cite{agrawal2023sarathi}.
Ultra-Nano-vLLM records prefill and decode time and throughput separately in
addition to end-to-end measurements. This separation supports phase-specific
profiling without treating a phase-local metric as a substitute for the total
batch result.

The default repository experiments use CUDA Graph execution, while an eager
mode is available as an explicit configuration. PyTorch documents CUDA Graph
replay as a way to reduce CPU launch overhead by replaying fixed kernels and
arguments~\cite{pytorchcudagraphs}. Because graph capture and its memory pools
are persistent runtime state, experiment configurations are isolated in fresh
processes.

\section{Method}

\subsection{Research Workflow}

Ultra-Nano-vLLM organizes performance work into five steps. A researcher first
establishes a reproducible scalar benchmark for the workload of interest. A
one- or two-dimensional experiment then isolates a workload or component
effect. The confirmed workload is profiled while its model, request count,
sequence lengths, sampling settings, execution mode, and component selectors
remain fixed. A focused component, kernel, or engine change is introduced only
after a bottleneck hypothesis has been formed. Finally, the original benchmark
grid is rerun and its reports are compared with the baseline.

\subsection{Scalar Benchmark Lifecycle}

The scalar benchmark configuration records the model path, request count,
primary and optional alternate input lengths, output length, maximum model
length, random seed, sampling temperature, repeat count, execution mode, and
engine-component selectors. An optional input-order field controls whether the
two input-length cohorts are alternating, primary-first, or alternate-first.
Both input lengths are checked against the requested and effective model
context limits before warmup.

Each benchmark creates the engine, constructs a deterministic synthetic
workload, performs warmup, and measures the requested repeats. CUDA is
synchronized around each end-to-end timing interval. The result contains
request and token totals; elapsed time; request-latency percentiles; request,
output-token, total-token, prefill, and decode throughput; prefill and decode
time; and peak KV-cache block use. Repeated scalar measurements are summarized
by their median, minimum, maximum, mean, and standard deviation and persisted
as typed JSON.

\subsection{Controlled Experiment Expansion}

An experiment YAML file defines workload columns and component columns.
Scalars and one-element lists broadcast within their respective set, while
other lists must have consistent lengths. The configuration layer normalizes
the input and validates labels, selector values, scalar configurations, and
the complete workload-by-component cross-product before model loading. The
sweep layer then zips positional columns within each axis and crosses the two
axes in deterministic workload-first order.

Each YAML file remains an independent reporting and plotting group. A group
may contain zero, one, or two conceptual dimensions. Zero-dimensional groups
produce the result and JSON report; one-dimensional groups produce line plots;
and two-dimensional groups produce multi-series line plots. Plot-only mode
loads only explicitly supplied reports and validates their compatibility
before rendering comparisons.

\subsection{Process Isolation and Reporting}

Every scalar experiment configuration runs sequentially in a fresh process
created with the spawn multiprocessing context. The worker performs benchmark
cleanup before returning the typed result and runtime configuration to the
parent. This boundary prevents CUDA Graph pools, cached allocations, and
distributed state from leaking into a later configuration while keeping file
output and plotting in the parent process.

The reporting layer produces thirteen comparisons: elapsed time; latency P50,
P90, and P99; request, output-token, total-token, prefill, and decode
throughput; prefill and decode time; peak KV-cache blocks; and peak KV-cache
utilization. For summarized metrics, plots use the mean with a
standard-deviation band. JSON reports retain the resolved experiment
configuration and runtime configuration required for later plot-only analysis.

\subsection{Pluggable Components and Kernels}

The \texttt{EngineComponent} value contains selectors for Scheduler,
BlockManager, Attention, Sampler, and the store-KV-cache kernel. Each selector
is a trusted local filename stem resolved only within the corresponding role
directory; path separators and parent-directory traversal are rejected.
Component files export \texttt{create\_component(...)} and kernel files export
\texttt{create\_kernel()}. Their return values must satisfy the
runtime-checkable protocol for the selected role.

This design makes implementation choice an experimental dimension. Benchmark
and experiment entrypoints transport selector values but do not import
concrete implementations or branch on component behavior. A comparison can
therefore replace one role while holding the workload and all other roles
fixed.

\subsection{Profiling Boundary}

Profiling begins only after an uninstrumented benchmark or controlled
experiment has identified a stable workload. The repository does not prescribe
a profiler. It provides optional NVTX ranges for scheduler operations, block
allocation and deallocation, model-input preparation, attention, KV-cache
storage, sampling, and logits computation. Profiling overhead is kept separate
from the JSON reports used for the performance comparison.

\section{Experiments}

\subsection{Experimental Setup}

We evaluate the baseline \texttt{scheduler} and two candidate implementations,
\texttt{scheduler-v1} and \texttt{scheduler-v2}. The evaluation uses a
Qwen3-0.6B model in bfloat16 on one NVIDIA GeForce RTX 3080 Laptop GPU, Python
3.12.3, PyTorch 2.8.0 with CUDA 12.8, and Triton 3.4.0. Tensor parallelism is
one. CUDA Graph execution is enabled, the model context limit is 8192 tokens,
and the KV cache contains 424 blocks of 256 tokens at a configured GPU-memory
utilization of 0.9.

Each workload contains 33 requests and generates one token per request with
temperature zero. The \texttt{hol} prefix identifies workloads designed to
expose head-of-line effects. \texttt{hol-short-only} is the homogeneous
control: all 33 prompts contain 256 tokens, so the queue contains no long
prompt. Each \texttt{hol-long-first-L} workload places 16 prompts of length
$L$ before 17 prompts of length 256. The six suffixes identify the exact long
prompt length: \texttt{hol-long-first-1024},
\texttt{hol-long-first-2048}, \texttt{hol-long-first-3072},
\texttt{hol-long-first-4096}, \texttt{hol-long-first-5120}, and
\texttt{hol-long-first-6144} use $L$ values of 1024, 2048, 3072, 4096, 5120,
and 6144 tokens, respectively. These workloads measure how long prompts at the
head of the queue delay the following short prompts.

The remaining two workloads hold the prompt lengths at 256 and 4096 tokens
while changing only their order. \texttt{hol-alternating-4096} alternates
between the two lengths, beginning with a 256-token prompt, for 17 short and
16 long prompts. \texttt{hol-short-first-4096} places all 17 short prompts
before the 16 long prompts, providing an ordering control in which the short
requests do not wait behind long prefills. Every scheduler--workload
coordinate has ten measured runs with seeds 1--10, for 270 measured batches
in total. The block manager, attention, sampler, and KV-cache storage kernel
remain fixed across scheduler variants.

\subsection{Scheduler Variants and Analysis}

The baseline scheduler processes the waiting queue in first-come,
first-served order and stops the prefill scan when the next sequence cannot
join the current batch. \texttt{scheduler-v1} performs a single pass over the
waiting queue, defers candidates that do not fit, and may backfill the batch
with a later request. \texttt{scheduler-v2} applies the same single-pass rule
after a stable sort by remaining prefill tokens, thereby implementing
shortest-remaining-job-first prefill scheduling. These policies leave the
decode batching logic unchanged, but both candidates select the running
sequence with the fewest allocated blocks as a preemption victim instead of
selecting the last running sequence. No preemption occurs in the measured
runs.

The analysis pairs each candidate run with the baseline run at the same
workload and seed. We report geometric-mean relative changes in
short-request median latency and total-token throughput. Confidence intervals
are the 2.5th and 97.5th percentiles of 10,000 paired bootstrap resamples.

\subsection{Performance Results}

Figure~\ref{fig:scheduler-improvement} shows that shortest-remaining-job-first
prefill scheduling substantially reduces head-of-line latency while keeping
aggregate throughput close to the baseline. Across the six long-first
workloads, \texttt{scheduler-v2} reduces short-request P50 by 25.15--91.66\%.
Its total-throughput change ranges from a 5.22\% reduction at 2048 tokens to a
1.16\% increase at 6144 tokens. The latency gain grows as the long prompt
becomes more expensive, reaching 83.51\%, 84.73\%, and 91.66\% for long prompts
of 4096, 5120, and 6144 tokens, respectively.

The alternating 4096-token workload provides the clearest result for both
candidates. \texttt{scheduler-v1} reduces short-request P50 by 61.58\% while
increasing throughput by 0.80\%; \texttt{scheduler-v2} reduces it by 71.56\%
with a 0.87\% throughput increase. In contrast, neither candidate materially
changes short-request latency in the homogeneous short-only workload or when
short prompts are placed first. This contrast localizes the benefit to
workloads in which short requests would otherwise wait behind long prefills.

\begin{figure*}[t]
\centering
\includegraphics[width=\textwidth]{experiments/report/scheduler-improvement.pdf}
\caption{The candidate schedulers reduce short-request head-of-line latency
with limited change in aggregate throughput. Panel (a) reports the reduction
in short-request P50 latency; panel (b) reports the change in total-token
throughput. Values are geometric-mean changes relative to the baseline, and
error bars are 95\% confidence intervals from 10,000 paired bootstrap
resamples over ten seeds. Positive values favor the candidate in both panels.
LF, Alt, and SF denote long-first, alternating, and short-first workloads; the
suffix gives the long-prompt length.}
\label{fig:scheduler-improvement}
\end{figure*}

\subsection{Mechanism Evidence and Limitations}

The telemetry is consistent with the intended scheduling changes, but it does
not establish causality. On the alternating 4096-token workload,
\texttt{scheduler-v1} records 120 backfill selections across ten runs
and raises mean prefill utilization from 0.711 to 0.853; the baseline records
no backfills. This change coincides with simultaneous latency and throughput
improvement on that workload. In contrast, \texttt{scheduler-v2} records no
backfills in any workload. It inspects more candidates than the baseline on
most heterogeneous workloads and its mean scheduler time is higher in those
cases, consistent with the added sorting and scan work.

The current experiment is bounded to one model, one GPU, 33-request offline
batches, one-token outputs, and nine prompt-order conditions. It therefore
tests prefill head-of-line behavior, not sustained decode scheduling or online
arrival processes. Experiments with longer outputs and KV-pressure workloads
are needed to determine whether the observed short-request latency reductions
persist outside this prefill-focused setting.

\section{Related Work}

LLM serving systems have studied scheduling, memory management, and the
interaction between prefill and decode. Orca introduced iteration-level
scheduling and selective batching for transformer-based generative
models~\cite{yu2022orca}. vLLM introduced PagedAttention and a block-based
KV-cache manager~\cite{kwon2023vllm}. Sarathi combined chunked prefills with
decode-maximal batching~\cite{agrawal2023sarathi}, while DistServe separated
prefill and decode execution to address interference between the two
phases~\cite{zhong2024distserve}. These systems propose serving mechanisms;
Ultra-Nano-vLLM instead supplies a small environment for isolating,
instrumenting, and comparing focused changes to the nano-vLLM codebase.

Reproducible benchmarking is a separate concern from proposing a serving
mechanism. MLPerf Inference defines a benchmarking methodology intended to
support comparison across inference systems~\cite{reddi2019mlperf}.
Ultra-Nano-vLLM has a narrower scope: it holds a local model and workload
configuration fixed while comparing components inside one lightweight engine.
Its reports and plots are research artifacts for controlled optimization, not
claims of conformance with MLPerf.

\section{Conclusion}

Ultra-Nano-vLLM extends nano-vLLM with a performance-research workflow that
connects scalar benchmarking, validated controlled experiments, process
isolation, typed reporting, plotting, profiling, and pluggable component
selection. The separation between engine-agnostic entrypoints and
file-selected implementations allows a focused component or kernel change to
be evaluated against the same workload and reporting path. The repository also
defines baseline and scheduler-comparison protocols that preserve per-repeat
measurements and correctness checks for later analysis.

\begin{thebibliography}{00}

\bibitem{kwon2023vllm}
W. Kwon et al., ``Efficient memory management for large language model serving
with PagedAttention,'' arXiv:2309.06180, 2023.

\bibitem{nanovllm}
GeeeekExplorer, ``nano-vLLM,'' GitHub repository. [Online]. Available:
\url{https://github.com/GeeeekExplorer/nano-vllm}. Accessed: Aug. 22, 2026.

\bibitem{reddi2019mlperf}
V. J. Reddi et al., ``MLPerf inference benchmark,'' arXiv:1911.02549, 2019.

\bibitem{agrawal2023sarathi}
A. Agrawal et al., ``SARATHI: Efficient LLM inference by piggybacking decodes
with chunked prefills,'' arXiv:2308.16369, 2023.

\bibitem{pytorchcudagraphs}
PyTorch Contributors, ``CUDA semantics: CUDA Graphs,'' PyTorch documentation.
[Online]. Available: \url{https://pytorch.org/docs/stable/notes/cuda.html}.
Accessed: Aug. 22, 2026.

\bibitem{yu2022orca}
G.-I. Yu, J. S. Jeong, G.-W. Kim, S. Kim, and B.-G. Chun, ``Orca: A
distributed serving system for transformer-based generative models,'' in
\emph{Proc. 16th USENIX Symp. Operating Systems Design and Implementation
(OSDI)}, 2022, pp. 521--538.

\bibitem{zhong2024distserve}
Y. Zhong et al., ``DistServe: Disaggregating prefill and decoding for
goodput-optimized large language model serving,'' arXiv:2401.09670, 2024.

\end{thebibliography}

\end{document}
